\scanpage{6}
\mspage{- 56 -}
n\textsuperscript{o}357
\[
0<\frac{1}{(f_{\unsure{o|0}}:f)}\mathrel{\handwritten{\leq}}\frac{(f:g)}{(f_o:g)}\mathrel{\handwritten{\leq}}(f:f_o)\ .
\]
\begingroup\csname tagsleft@true\endcsname\[
(f+f':g)\mathrel{\handwritten{\leq}}(f:g)+(f':g).\tag*{(vii)}
\]\endgroup

(viii) \quad si $f$, $f'$, $h$ sont dans $\handwritten{\unsure{\mathcal{K}|\mathscr{K}}}_+$, si $h\mathrel{\handwritten{\geq}}1$ sur le support de $f+f'$, et si $\varepsilon>0$, il y a un voisinage compact $V$ de $e$ tel que l'on ait, pour tout $g\mathrel{\handwritten{\in}}\handwritten{\mathcal{K}}^{\ast}_{+}(V)$ :
\[
(f+f':g)\mathrel{\handwritten{\geq}}(f:g)+(f':g)-\varepsilon(h:g)\ .
\]

Les trois premières sont évidentes ; (iv) aussi, puisque $\sup(\handwritten{\gamma}(\mu)g)\mathrel{\handwritten{\leq}}\lVert\mu\rVert\sup(g)$ ; (v) résulte aussitôt de la formule $\lVert\mu\ast\mu'\rVert={ }=\lVert\mu\rVert.\lVert\mu'\rVert$ (par la vertu du Chap. VII, \handwritten{\S}\ 1, n\textsuperscript{o}1, ou, moins savamment et plus simplement dans le cas qui nous occupe ici, par un calcul direct trivial) et de ce qu'on a, si $f\mathrel{\handwritten{\leq}}\handwritten{\gamma}(\mu)g$ et $g\mathrel{\handwritten{\leq}}\handwritten{\gamma}(\mu')h$ :
\[
f\mathrel{\handwritten{\leq}}\handwritten{\gamma}(\mu)\handwritten{\gamma}(\mu')h=\handwritten{\gamma}(\mu\ast\mu')h\ .
\]

Si on applique (v) à $f_o$, $f$, $g$ d'une part, et à $f$, $f_o$, $g$ d'autre part, on obtient (vi) ; (vii) est évident. Enfin, soient $f$, $f'$, $h$ comme il est dit dans (viii). Posons $F=f+f'+\handwritten{\delta} h$, avec $\handwritten{\delta}>0$ ; les fonctions $\varphi$, $\varphi'$ qui coïncident respectivement avec $f/F$ et $f'/F$ sur le support de $f+f'$, et qui sont $0$ en dehors de celui-ci, sont dans $\handwritten{\mathcal{K}}$, donc uniformément continues ; pour tout $\handwritten{\eta}>0$, on peut donc choisir un voisinage compact $V$ de $e$ tel que $|\varphi(x)-\varphi(s)|$ et $|\varphi'(x)-\varphi'(s)|$ soient $\mathrel{\handwritten{\leq}}\handwritten{\eta}$ pour $s^{-1}x\mathrel{\handwritten{\in}}V$ ; alors, quelle que soit $g\mathrel{\handwritten{\in}}\handwritten{\mathcal{K}}^{\ast}_{+}(V)$, on aura
\[
\varphi(x)g(s^{-1}x)\mathrel{\handwritten{\leq}}\unsure{[|(}\varphi(s)+\handwritten{\eta}]g(s^{-1}x)
\]
quels que soient $x$ et $s$, et de même pour $\varphi'$. Supposons maintenant qu'on ait, pour tout $x$ :
\[
F(x)\mathrel{\handwritten{\leq}}\handwritten{\sum}_{i=1}^{n}a_i g(s_i^{-1}x)\ .
\]
Puisque $f=\varphi F$, on en conclut :
% archive sheet number: 5

